% !TeX root = ../../Book2.tex
% 纲要标题：### 17.2 中心化子、子域与分裂域
% 纲要位置：计划纲要.md，第 1313 行。

\section{中心化子、子域与分裂域}
\label{b2:sec:1702}

一个单子代数 \(B\subseteq A\) 的中心化子记录与它交换的另一部分乘法。维数公式将精确说明这两部分各占多大；但只有 \(B\) 的中心也是 \(k\) 时，二者才直接张量成整个 \(A\)。本节通过简单模及其重数证明这些结论，子代数中心的扩张 \(Z(B)/k\) 可以不可分。

% 纲要要点（供目录核对）：
%
% * 中心化子的维数关系
% * 双中心化子定理
% * 子域与最大交换子代数
% * 可分最大子域与有限 Galois 分裂域
%

\subsection{中心化子的维数关系}
\label{b2:subsec:1702:01}

记
\[
C_A(B)=\{\,a\in A\mid ab=ba\text{ 对所有 }b\in B\,\}.
\]
它是含单位元的 \(k\) 子代数。本节的中心化子在 \(A\) 内取，与第十一章在整个线性自同态代数内取中心化子应分别理解。
\symidx{\(C_A(B)\)：子代数 \(B\) 在 \(A\) 中的中心化子}

\begin{theorem}[中心化子定理]\label{b2:thm:csa-centralizer}
设 \(k\) 为域，\(A\) 为中心单 \(k\) 代数，\(B\subseteq A\) 为非零简单含幺 \(k\) 子代数，\(L=Z(B)\)。令 \(C=C_A(B)=\{\,a\in A\mid ab=ba\text{ 对所有 }b\in B\,\}\)，则
\[
C\text{ 为简单代数},\qquad Z(C)=L,
\qquad \dim_kA=\dim_kB\,\dim_kC,
\qquad C_A(C)=B.
\]
其中两个中心均按它们在 \(A\) 中的实际嵌入识别。结论不要求 \(L/k\) 可分。
\end{theorem}
\begin{proof}
先构造有限维代数 \(E=B\otimes_kA^{\mathrm{op}}\)。\(B\) 以其中心 \(L\) 为底域是中心单代数，而
\[
E\cong B\otimes_L(L\otimes_kA^{\mathrm{op}}).
\]
具体同构将 \(b\otimes(l\otimes a^{\mathrm{op}})\) 送到 \(bl\otimes a^{\mathrm{op}}\)，逆映射在 \(b\otimes a^{\mathrm{op}}\) 上取 \(b\otimes(1\otimes a^{\mathrm{op}})\)。\cref{b2:thm:csa-base-change}说明第二个因子在 \(L\) 上中心单，再由\cref{b2:thm:csa-tensor-product}，\(E\) 为中心单 \(L\) 代数，特别是单 Artin 环。

令 \(E\) 在 \(A\) 上作用为 \((b\otimes a^{\mathrm{op}})x=bxa\)。一个 \(E\) 线性自同态首先与全部右乘交换，所以由\cref{b2:thm:regular-double-centralizer}，它形如左乘 \(c\in A\)。再要求它与左 \(B\) 作用交换，恰好得到 \(c\in C\)。左乘保留乘法次序，因此
\[
C\cong\operatorname{End}_E(A).
\]
写 \(E\cong M_r(D)\)，其中 \(D\) 是除环。\cref{b2:thm:artin-wedderburn}及\cref{b2:prop:semisimple-multiplicity}的矩阵块分类说明，\(E\) 半单，且简单左模只有列模 \(D^r\) 一种类型。\(A\) 是有限维 \(k\) 空间，所以作为 \(E\) 模为有限个副本之和，设 \(A\cong(D^r)^m\)、\(m\ge1\)。因此
\[
\operatorname{End}_E(A)\cong M_m(D^{\mathrm{op}}),
\]
从而 \(C\) 简单。

记 \(d=\dim_kD\)。分别比较上述代数与模的 \(k\) 维数，有
\[
\dim_kB\,\dim_kA=r^2d,\qquad
\dim_kA=mrd,\qquad\dim_kC=m^2d.
\]
这三式给出 \(\dim_kB\,\dim_kC=\dim_kA\)。至此已独立证明：对任意满足定理假设的简单子代数，其中心化子简单，且有上述维数公式；双中心化子和中心的识别尚未用到。现在 \(C\) 也是 \(A\) 的简单含幺子代数，将这个已证的中间结论用于 \(C\subseteq A\)，得到
\(\dim_k C_A(C)=\dim_kA/\dim_kC=\dim_kB\)。而 \(B\subseteq C_A(C)\) 由交换性定义直接成立，有限维与等维使二者相等。最后
\[
Z(C)=C\cap C_A(C)=C\cap B=Z(B)=L,
\]
得到中心的实际识别。
\end{proof}

证明中 \(L\otimes_kA^{\mathrm{op}}\) 的中心单性来自对任意域扩张均成立的标量扩张定理。即使 \(L/k\) 是纯不可分扩张，这一步也仍然有效；真正需要简单性的，是 \(B\) 以其中心为底域的结构。

\subsection{双中心化子定理}
\label{b2:subsec:1702:02}

\begin{corollary}\label{b2:cor:csa-centralizer-factorization}
设 \(k\) 为域，\(A\) 为中心单 \(k\) 代数，\(B\subseteq A\) 为非零简单含幺 \(k\) 子代数，记 \(C_A(B)\) 为 \(B\) 在 \(A\) 中的中心化子。若 \(Z(B)=k\)，则乘法诱导中心单 \(k\) 代数同构
\[
B\otimes_kC_A(B)\longrightarrow A,\qquad b\otimes c\longmapsto bc.
\]
一般地，令 \(L=Z(B)\)，则有 \(L\) 代数同构
\[
B\otimes_LC_A(B)\cong C_A(L).
\]
\end{corollary}
\begin{proof}
两个因子的元素互相交换，所以乘法公式保持张量关系与代数乘法。若 \(L=k\)，来源由\cref{b2:thm:csa-tensor-product}为单代数；映射保持单位元，故核不能为全代数，只能为零。中心化子维数公式说明来源与目标同维，于是满射。

一般情形中，\(B,C_A(B)\) 都是中心单 \(L\) 代数，故它们在 \(L\) 上的张量积单，乘法仍单射，且其像包含于 \(C_A(L)\)。这里仍以 \(k\) 为底域，将简单 \(k\) 子代数 \(L\subseteq A\) 代入中心化子定理，得到 \(\dim_k C_A(L)=\dim_kA/[L:k]\)。张量积来源的 \(k\) 维数为
\(\dim_kB\,\dim_kC_A(B)/[L:k]\)，也等于这个数，因此像就是全部 \(C_A(L)\)。
\end{proof}

例如，把有限扩张 \(L/k\) 通过正则作用嵌入 \(A=\operatorname{End}_k(L)\)。与左乘 \(L\) 交换的映射恰为右乘 \(L\)，因为 \(T(x)=T(x\cdot1)=x\cdot T(1)\)。\(L\) 交换，所以 \(C_A(L)=L\)，维数公式给出 \([L:k]^2=[L:k]\,[L:k]\)。当 \(L\ne k\) 时，\(L\otimes_L L=L\) 远小于 \(A\)；只有在正确的中心上取张量，且把目标写成 \(C_A(L)\)，才能保持一般公式成立。

\ProseParagraphBreak
双中心化子等式还说明，简单子代数可由它在 \(A\) 中的全部交换关系恢复。若删除简单性，这一性质可能失败：第十一章的上三角代数 \(T_2(k)\subseteq M_2(k)\) 的中心化子为 \(kI\)，再取一次得到整个 \(M_2(k)\)，并非原上三角代数。

这些关系可以在 \(A\) 内按包含画成两支：
\[
\begin{tikzcd}[column sep=2.2em,row sep=1.5em]
& & B \arrow[dr,hook] & & \\
k \arrow[r,hook] & L=Z(B)=Z(C) \arrow[ur,hook] \arrow[dr,hook]
& & C_A(L) \arrow[r,hook] & A\\
& & C=C_A(B) \arrow[ur,hook] & &
\end{tikzcd}
\]
两支相乘得到 \(B\otimes_LC\cong C_A(L)\)，而 \(C_A(C)=B\)、\(\dim_kA=\dim_kB\,\dim_kC\)。当 \(L=k\) 时，图中的 \(C_A(L)\) 才是整个 \(A\)。

\subsection{子域与最大交换子代数}
\label{b2:subsec:1702:03}

设 \(A\) 为中心单 \(k\) 代数。若含 \(k\) 的子域 \(L\subseteq A\) 不包含于更大的 \(k\) 子域，称 \(L\) 为\term[baohan-jida-ziyu]{包含极大子域}[inclusion-maximal subfield]。若一个含幺交换子代数不包含于更大的含幺交换子代数，则称它为\term[zuida-jiaohuan-zidaishu]{最大交换子代数}[maximal commutative subalgebra]。前者只在域中比较，后者也允许零因子或幂零元。文献中的 maximal subfield 还可能专指满足 \([L:k]=\deg A\) 的子域；本书在一般中心单代数中分别写明这两个条件，只在除代数语境把包含极大子域简称\term[zuida-ziyu]{最大子域}[maximal subfield]。

\cref{b2:tab:csa-subfield-conditions}按对象所在位置列出容易混淆的条件。最后一行的 \(K\) 只须是 \(k\) 的扩域，可以完全不嵌入 \(A\)。
\begin{table}[htbp]
\centering
\caption{中心单代数中的子域、交换子代数与分裂域}\label{b2:tab:csa-subfield-conditions}
\begin{tabularx}{\linewidth}{@{}p{0.31\linewidth}X@{}}
\toprule
对象 & 条件 \\
\midrule
包含极大子域 \(L\subseteq A\) & 不包含于 \(A\) 中更大的含 \(k\) 子域。\\
最大交换子代数 \(E\subseteq A\) & 不包含于 \(A\) 中更大的含幺交换 \(k\) 子代数。\\
自中心化子域 \(L\subseteq A\) & \(C_A(L)=L\)。\\
次数达到 \(\deg A=n\) 的子域 \(L\subseteq A\) & \([L:k]=n\)。\\
分裂域 \(K/k\) & \(A_K\cong M_n(K)\)；不要求 \(K\subseteq A\)。\\
\bottomrule
\end{tabularx}
\end{table}

\begin{proposition}\label{b2:prop:csa-self-centralizing-field}
设 \(k\) 为域，\(A\) 为中心单 \(k\) 代数，\(n=\deg A\)，\(L\subseteq A\) 为含 \(k\) 的子域，则以下条件等价：
\[
[L:k]=n;\qquad C_A(L)=L;\qquad
L\text{ 是 }A\text{ 的最大交换子代数}.
\]
若 \(A\) 本身为除代数，则每个含 \(k\) 的最大子域都满足这些条件。
\end{proposition}
\begin{proof}
将\cref{b2:thm:csa-centralizer}用于简单子代数 \(L\)，得到 \(n^2=[L:k]\dim_k C_A(L)\)。由于 \(L\subseteq C_A(L)\)，当 \([L:k]=n\) 时二者同维，所以相等；反向相等也立即给出次数条件。

若 \(C_A(L)=L\)，任何包含 \(L\) 的交换子代数都包含于该中心化子，因此只能是 \(L\)。反过来，若 \(c\in C_A(L)\)，则 \(L[c]\) 是包含 \(L\) 的交换子代数，最大交换性使 \(c\in L\)。最后，若环境为除代数 \(D\)，则 \(L[c]\) 是有限维交换整环，因而为域：非零元的乘法单射，有限维使其满射，从而有逆。最大子域的性质就迫使 \(L[c]=L\)，也得到中心化子等于 \(L\)。
\end{proof}

在矩阵代数中，包含极大不保证次数达到 \(\deg A\)。例如 \(M_2(\mathbb C)\) 中每个有限维的 \(\mathbb C\) 子域都等于 \(\mathbb C\)，所以标量域是包含极大子域，却不是最大交换子代数，也只有次数一。对角代数 \(\mathbb C\times\mathbb C\) 是最大交换子代数，但不是域；\(\mathbb C[N]\)、\(N=\begin{pmatrix}0&1\\0&0\end{pmatrix}\)，也是最大交换子代数，因为与 \(N\) 交换的矩阵恰为 \(aI+bN\)，而它含有非零幂零元。

\begin{proposition}\label{b2:prop:maximal-field-splits-csa}
设 \(k\) 为域，\(A\) 为中心单 \(k\) 代数，\(L\subseteq A\) 为含 \(k\) 的子域。若 \([L:k]=\deg A=n\)，则 \(L\) 分裂 \(A\)。特别地，每个有限维中心除 \(k\) 代数的最大子域都是它的分裂域。
\end{proposition}
\begin{proof}
把 \(A\) 看成右 \(L\) 向量空间，其维数为 \(\dim_kA/[L:k]=n\)。作用
\[
L\otimes_kA\longrightarrow\operatorname{End}_L(A),
\qquad l\otimes a\longmapsto(x\mapsto axl)
\]
是 \(L\) 代数同态；右 \(L\) 线性用到 \(L\) 内元素交换，乘法相容性由结合律与相同交换性给出。目标中的 \(A\) 始终是原来的右 \(L\) 空间。来源由\cref{b2:thm:csa-base-change}为单代数，所以保持单位元的映射单射。两侧的 \(L\) 维数均为 \(n^2\)，故为同构，目标为 \(M_n(L)\)。除代数的最大子域由\cref{b2:prop:csa-self-centralizing-field}满足 \([L:k]=\deg A\)，所以也适用。
\end{proof}

\subsection{可分最大子域与有限 Galois 分裂域}
\label{b2:subsec:1702:04}

有限维除代数确有最大子域：子域维数是有界正整数，取具有最大维数的一个即可。不过，这样得到的子域未必可分。为得到有限 Galois 分裂域，下面另作一次可分元素的选择。

\begin{theorem}[有限可分分裂域的存在]\label{b2:thm:csa-separable-splitting}
设 \(k\) 为域。每个中心单 \(k\) 代数都有有限可分分裂域，因而也有有限 Galois 分裂域。更具体地，每个有限维中心除 \(k\) 代数 \(D\) 都含有一个可分最大子域 \(L\)，满足 \([L:k]=\deg D\)，且 \(L\) 分裂 \(D\)。
\end{theorem}
\begin{proof}
由\cref{b2:prop:csa-division-description}，只须对中心除代数 \(D\) 找到所述 \(L\)：一旦 \(L\) 分裂 \(D\)，它也分裂 \(M_r(D)\)。若 \(k\) 有限，则 \(D\) 的元素总数也有限，由\cref{b2:thm:wedderburn-little}，\(D\) 交换；中心为 \(k\) 便迫使 \(D=k\)，取 \(L=k\) 即可。以下设 \(k\) 无限，令 \(n=\deg D\)，\(\Omega\) 为代数闭包。

取同构 \(D_\Omega\cong M_n(\Omega)\) 及 \(D\) 的 \(k\) 基 \(e_1,\ldots,e_{n^2}\)，其像为矩阵 \(E_1,\ldots,E_{n^2}\)。这些矩阵是 \(M_n(\Omega)\) 的 \(\Omega\) 基。要寻找 \(a=\sum_i c_ie_i\)、\(c_i\in k\)，使对应矩阵具有 \(n\) 个互异特征值。

先建立一个明确的多项式检验。对任意 \(H\in M_n(\Omega)\)，令 \(f_H(t)=\det(tI-H)\)，并令 \(U_j\) 表示次数小于 \(j\) 的 \(\Omega\) 多项式空间，\(U_0=0\)。考虑等维空间之间的映射
\[
U_{n-1}\oplus U_n\longrightarrow U_{2n-1},
\qquad(u,v)\longmapsto u f_H+v f_H'.
\]
在单项式基下取其行列式 \(\Delta(H)\)。这是首一多项式 \(f_H\) 与 \(f_H'\) 的结式 \(\operatorname{Res}(f_H,f_H')\) 的一种归一化，也可视为判别式；更换这些空间的固定基只会使行列式乘以一个非零常数。矩阵条目来自 \(f_H\) 及其导数的系数，所以 \(\Delta(H)\) 是 \(H\) 条目的多项式。

该行列式非零当且仅当 \(f_H,f_H'\) 互素。如果互素且 \(u f_H+v f_H'=0\)，则 \(f_H\mid v\)，但 \(\deg v<n\)，故 \(v=0\)，继而 \(u=0\)，映射单射而可逆。若有正次数公因子 \(h\)，则 \((u,v)=(f_H'/h,-f_H/h)\) 是符合次数限制的非零核向量，所以行列式为零。互素又等价于 \(f_H\) 没有重根：根 \(\lambda\) 同时使 \(f_H,f_H'\) 为零，恰好是它在因式分解中出现至少两次。

于是
\[
P(X_1,\ldots,X_{n^2})=\Delta\left(\sum_iX_iE_i\right)
\in\Omega[X_1,\ldots,X_{n^2}]
\]
是非零多项式。因为 \(E_i\) 张成全部矩阵，可在 \(\Omega\) 中选择系数使其和为具有 \(n\) 个互异对角元的对角矩阵，此处 \(\Delta\ne0\)。

即使 \(P\) 的系数在 \(\Omega\) 中，它也不能在无限子域 \(k\) 的全部点上消失。证明可对变量数归纳：一元非零多项式只有有限多个根；多元时先将它按最后一个变量展开，选一个非零系数多项式，用归纳假设为前面的变量取 \(k\) 中的值，使该系数非零，再用一元结论选最后一个值。因此可取 \(c_i\in k\) 使 \(P(c_1,\ldots,c_{n^2})\ne0\)。相应的 \(a=\sum_i c_ie_i\in D\) 在分裂矩阵模型中有 \(n\) 个互异特征值。任意零化矩阵的多项式必须在这 \(n\) 个特征值处为零，故次数至少为 \(n\)；其 \(n\) 次特征多项式又由\cref{b2:thm:exterior-cayley-hamilton}零化该矩阵。因此矩阵最小多项式就是这个没有重根的特征多项式。

\(a\) 在 \(k\) 上的最小多项式也正是上述矩阵的最小多项式。事实上，对左乘算子 \(L_a:D\to D\)，有 \(p(L_a)=0\) 当且仅当 \(p(a)=0\)：必要性只需将算子作用于 \(1\)。\cref{b2:prop:scalar-extension-invariant-factors}说明 \(L_a\) 扩张到 \(\Omega\) 后最小多项式原样不变；在 \(D_\Omega\cong M_n(\Omega)\) 上，这个扩张算子就是矩阵 \(H\) 的左乘，其最小多项式又等于 \(H\) 的最小多项式。因此 \(a\) 的首一最小多项式次数为 \(n\)，且在 \(\Omega\) 中无重根，所以这个最小多项式可分。这里没有预先假定 \(H\) 的特征多项式系数属于 \(k\)。

\(D\) 为除环，使交换子代数 \(k[a]\) 是有限维整环，因而为域。于是 \(L=k[a]\) 是次数 \(n\) 的可分子域，由\cref{b2:prop:csa-self-centralizing-field}及\cref{b2:prop:maximal-field-splits-csa}，它是最大子域并分裂 \(D\)。最后把 \(a\) 送到其最小多项式在 \(\Omega\) 中的一个根，将 \(L\) 嵌入 \(\Omega\)，并取该多项式在 \(\Omega\) 中的分裂域 \(F\)。第一卷命题15.1.2保证其有限性；它由这个可分多项式的全部根生成，按该卷定义19.1.1后的分裂域刻画，是有限 Galois 扩张，并且含有所选的 \(L\) 的像。继续扩张分裂矩阵环仍为矩阵环，所以 \(F\) 也分裂原来的 \(A\)。
\end{proof}
